\section{近端与坐标方法利用可分结构}
\label{sec:09-Convex-Optimization/06-Optimization-Algorithms/04}

接着上一篇那份 \(\ell_1\) 正则的模型往下做。特征有两万列，交付要求是一张不超过五十行的特征清单，而次梯度法跑了一万步之后，最小的那批系数停在 \(10^{-7}\) 附近，一个精确的零都没有。手工加一道阈值把小系数抹掉当然可以，但阈值取多少没有依据，抹完之后的模型也不再是任何优化问题的解——它只是一个被截断过的近似。

回头看问题出在哪里。目标由两块拼成：

\[
F(w)=\underbrace{\frac{1}{2}\norm{Xw-y}_2^2}_{f,\ \text{光滑}}
+\underbrace{\lambda\norm{w}_1}_{g,\ \text{非光滑但形状已知}} .
\]

次梯度法对这两块一视同仁，各取一个次梯度，加起来当成一个方向。这么做等于把 \(g\) 也当成只能查询斜率的黑箱，而 \(\lambda\norm{w}_1\) 的全部信息明明写在纸上，逐坐标可分，一维子问题能手算。信息被自己丢掉了，慢和不稀疏都是这次丢弃的后果。

\subsection{把两块用各自合适的信息处理}

第一篇里的下降引理给了光滑部分一个可用的局部模型：在 \(x_k\) 处线性化 \(f\)，再补一个由 \(L\) 控制的二次项，就得到一条压在 \(f\) 上方、处处不被顶破的抛物线。非光滑的 \(g\) 不需要任何模型，它自己就足够简单。把两者拼在一起，下一步取

\[
x_{k+1}=\argmin_x\Bigl\{
f(x_k)+\nabla f(x_k)^{\trans}(x-x_k)+\frac{1}{2\alpha}\norm{x-x_k}^2+g(x)
\Bigr\}.
\]

前三项配方成一个完全平方，与 \(x\) 无关的部分丢掉，右端就化为

\[
x_{k+1}=\argmin_x\Bigl\{g(x)+\frac{1}{2\alpha}\norm{x-v_k}^2\Bigr\},
\qquad
v_k=x_k-\alpha\nabla f(x_k).
\]

括号里的运算对任何凸函数 \(g\) 都有意义，记作近端算子

\[
\operatorname{prox}_{\alpha g}(v)=\argmin_x\Bigl\{g(x)+\frac{1}{2\alpha}\norm{x-v}^2\Bigr\}.
\]

二次项强凸，所以只要 \(g\) 是闭凸的，这个极小点存在且唯一。整个迭代读作：光滑部分走一步梯度，得到 \(v_k\)；非光滑部分做一次近端修正，得到 \(x_{k+1}\)。

值得停一下的是这两步在信息使用上的不对称。\(f\) 被线性化了，由此产生的近似误差由 \(L\) 兜底；\(g\) 一个字都没近似，它以完整形式进入子问题。次梯度法把 \(g\) 压成一条支撑直线，那条直线在尖角处可以任意倾斜，误差不受控；近端步保留了尖角本身，于是尖角带来的那个特殊解——精确的零——才有机会被取到。

\subsection{\texorpdfstring{软阈值是 \(\ell_1\) 的近端算子}{软阈值是 L1 的近端算子}}

对 \(g=\lambda\norm{\cdot}_1\)，近端子问题按坐标完全解耦，一维的问题是

\[
\min_x\ \lambda\abs{x}+\frac{1}{2\alpha}(x-v)^2 .
\]

用上一篇的最优性条件 \(0\in\partial(\cdot)\)。\(x>0\) 时次微分是单点，条件为 \(\lambda+(x-v)/\alpha=0\)，解得 \(x=v-\alpha\lambda\)，与假设相容当且仅当 \(v>\alpha\lambda\)；\(x<0\) 对称地给出 \(x=v+\alpha\lambda\)，要求 \(v<-\alpha\lambda\)；\(x=0\) 时次微分是区间 \([-\lambda,\lambda]\)，条件成为 \(-v/\alpha\in[-\lambda,\lambda]\)，即 \(\abs{v}\le\alpha\lambda\)。三种情形覆盖了整条实线且互不重叠，合起来就是软阈值算子

\[
S_{\alpha\lambda}(v)=
\begin{cases}
v-\alpha\lambda,&v>\alpha\lambda,\\
0,&\abs{v}\le\alpha\lambda,\\
v+\alpha\lambda,&v<-\alpha\lambda,
\end{cases}
\]

向量情形逐坐标套用。一行闭式解，成本与坐标数成正比。

\begin{center}
\begin{tikzpicture}[x=1.15cm,y=1.0cm,>=stealth]
  \draw[->,black!40] (-2.65,0) -- (2.8,0) node[right,font=\footnotesize] {\(v\)};
  \draw[->,black!40] (0,-2.5) -- (0,2.55);
  \draw[dashed,black!45] (-2.3,-2.3) -- (2.3,2.3);
  \node[font=\footnotesize,black!50] at (2.72,2.25) {\(y=v\)};
  \draw[black!85,very thick] (-2.3,-1.7) -- (-0.6,0) -- (0.6,0) -- (2.3,1.7);
  \fill[black] (-0.6,0) circle (1.5pt);
  \fill[black] (0.6,0) circle (1.5pt);
  \node[below,font=\footnotesize] at (-0.62,-0.06) {\(-\tau\)};
  \node[below,font=\footnotesize] at (0.64,-0.06) {\(\tau\)};
  \node[font=\footnotesize,black!70] at (-1.62,1.0) {阈值内整段压成零};
  \draw[->,black!55] (-1.62,0.85) -- (-0.2,0.06);
  \draw[<->,black!55] (1.6,1.6) -- (1.6,1.0);
  \node[right,font=\footnotesize,black!55] at (2.45,1.3) {外侧平移 \(\tau\)};
  \draw[black!55] (1.66,1.3) -- (2.4,1.3);
  \node[above left,font=\footnotesize,black!70] at (-2.3,-1.68) {\(S_\tau(v)\)};
\end{tikzpicture}

\emph{软阈值把 \([-\tau,\tau]\) 整段折到零，区间外的部分保留斜率 \(1\) 但向原点平移 \(\tau\)；稀疏来自那段被压平的死区，收缩来自那个平移。}
\end{center}

图里两件事各有出处。中间那段平的来自 \(x=0\) 处次微分是一个区间：只要 \(v\) 落进 \([-\tau,\tau]\)，正则项的``摩擦''就压得住数据项的推力，最优解精确取零，而不是取一个很小的数。两侧那个平移量 \(\tau=\alpha\lambda\) 则是 \(\ell_1\) 对留下来的系数收的租金——这正是 Lasso 的估计相对最小二乘存在收缩偏差的原因，见\hyperref[sec:13-Machine-Learning/02-Linear-Models/02]{``正则化从 Ridge 到 Lasso''}。同一件事在几何上的画法，即 \(\ell_1\) 球的尖角为何顶在坐标轴上，见\hyperref[sec:09-Convex-Optimization/07-Applications/02]{``\(L_1\) 正则为何产生稀疏''}。

对照一下深度模型里常见的稀疏化做法就更清楚。训练完再按幅值裁掉小权重，是先解一个完全没有稀疏偏好的问题，再手工截断，截断点是外生的，截完的模型不满足任何最优性条件；而把 \(\ell_1\) 写进目标、用近端步训练，零是解自己产生的，阈值由 \(\lambda\) 和步长共同决定，剪掉的坐标满足上一篇那个不等式 \(\abs{\partial L/\partial w_j}\le\lambda\)。两种做法都能给出稀疏的权重，但只有后者能说清楚被剪掉的东西为什么可以剪。

\subsection{投影梯度是同一个公式的特例}

近端算子还覆盖了带约束的情形。取 \(g\) 为凸集 \(C\) 的指示函数——在 \(C\) 内取 \(0\)，在 \(C\) 外取 \(+\infty\)——近端子问题变成在 \(C\) 内找离 \(v\) 最近的点：

\[
\operatorname{prox}_{\alpha I_C}(v)=\argmin_{x\in C}\norm{x-v}^2=\Pi_C(v),
\]

于是迭代化为投影梯度法 \(x_{k+1}=\Pi_C\bigl(x_k-\alpha\nabla f(x_k)\bigr)\)。约束和正则在这个框架里是同一件事的两种写法，这也解释了上一篇那个投影次梯度步为什么能沿用原来的证明：投影不增加到可行点的距离。至于投影本身的几何，见\hyperref[sec:03-Linear-Algebra/03-Orthogonality/02]{``投影与最近点：从一条直线走向一般子空间''}。

这个特例同时划出了方法的适用边界。盒子、球、单纯形、非负象限的投影都有闭式解或近似闭式的算法，代价与维数同阶。可一旦 \(C\) 由一组一般的不等式围成，投影本身就是一个和原问题同样难的优化问题，见\hyperref[sec:09-Convex-Optimization/02-Convex-Sets/02]{``等式与不等式如何塑造可行域''}。近端方法要求那个子问题便宜，这不是技术细节，而是它能不能用的前提。

\subsection{速率回到光滑情形的水平}

这一步换回了什么，看速率。设 \(\nabla f\) 是 \(L\)-Lipschitz，取 \(\alpha\le1/L\)，近端梯度法对复合目标给出

\[
F(x_k)-F^\star=O(1/k),
\]

和第一篇里纯光滑问题的速率相同。非光滑项没有拖慢任何东西，因为它被精确求解，不产生近似误差；决定速率的只剩 \(f\) 的光滑常数。上一篇那个 \(O(1/\sqrt k)\) 的下界并没有被推翻——它约束的是只能查询次梯度的算法，而近端步用了远多于一个次梯度的信息。

Nesterov 型外推可以照搬过来，在两次近端步之间插入一个动量项，就是 FISTA，速率提到 \(O(1/k^2)\)。同样的注意事项也照搬：加速后的目标值不再单调，非光滑处附近可能出现振荡，实践中常用单调变体或自适应重启把速率换成稳定。另一个容易忽略的收益是稀疏模式会稳定下来——迭代几十步之后，取零的坐标集合基本不再变化，后续的矩阵向量乘法只需在活跃坐标上做，每步成本随之下降。

\subsection{坐标下降把可分性用到底}

近端方法利用的是加法可分：\(g\) 拆成各坐标的和。若这种可分性还延伸到子问题层面，可以更激进——一次只动一个坐标，其余固定。对 Lasso，固定 \(w_{-j}\) 之后关于 \(w_j\) 的子问题是一个一维的最小二乘加绝对值，解正是上面那个软阈值，一行算完。一轮扫过所有坐标的成本主要花在残差更新上，而设计矩阵稀疏时这个更新只触及少数几行，这就是 \texttt{glmnet} 一类求解器能在几万维上秒级出解的原因。

坐标的选取可以循环，可以随机，也可以按当前的坐标梯度大小挑。收敛的前提是非光滑部分可分：若 \(g\) 里出现 \(\norm{w}_1\) 这样的逐坐标结构，坐标下降收敛到最优；若换成组稀疏或一个把坐标耦合起来的非光滑项，只动一个坐标可能在某个非最优点上卡死——沿任何单坐标方向都无法改进，联合移动却可以。此时要按``块''来更新，块的划分要把耦合的坐标放在一起。交替最小化的这套思路在别处也反复出现，例如\hyperref[sec:13-Machine-Learning/05-Clustering-and-Data-Geometry/01]{``K-means 在交替下降中制造簇''}里对簇分配和簇中心的交替，只是那里的目标不再凸。

选哪种方法，看的不是名字而是可用的计算原语。全梯度便宜、目标光滑，就用梯度法或它的加速版；非光滑项的 prox 有闭式解，就用近端梯度；变量高度可分、数据稀疏、单坐标更新闭式，坐标下降往往最快；什么都拿不到、只有一个次梯度接口，才退回上一篇的方法。

这条线索到这里也走到了它的边界。近端与坐标方法要求非光滑部分简单到能被精确求解，或者可行域简单到能被廉价投影。可现实中的约束经常是一组互相牵扯的不等式，投影不便宜，可分性也无从谈起，而我们仍然想要一个能给出对偶证书的高精度解。\hyperref[sec:09-Convex-Optimization/06-Optimization-Algorithms/05]{``从障碍函数到原始---对偶内点法''}换一条路：不去投影，而是把约束改写进目标，让迭代自己待在可行域里面。
